% Conceptual adaptations of ME 537 HW02 P1, P3--P6, P9; no textbook assets copied.
% Shared layout adapted from the instructor's ME 537 homework sources.
\documentclass[12pt,letterpaper]{article}
\usepackage{mathptmx}
\usepackage[T1]{fontenc}
\usepackage{xcolor,amsmath,amssymb,xurl,listings,enumitem,tikz}
\textwidth=6.5in
\textheight=9in
\oddsidemargin=0in
\evensidemargin=0in
\topmargin=0in
\headheight=0in
\headsep=0in
\parskip=0.08in
\parindent=0in
\setlist[enumerate]{itemsep=0.04in,topsep=0.04in}
\lstset{basicstyle=\ttfamily\footnotesize,breaklines=true,columns=fullflexible,
keepspaces=true,showstringspaces=false,aboveskip=0.04in,belowskip=0.04in}
\newcommand{\R}[2]{{}^{#1}\!R_{#2}}
\newcommand{\T}[2]{{}^{#1}\!T_{#2}}
\newcommand{\courseheader}[2]{%
\setlength{\fboxrule}{2pt}\setlength{\fboxsep}{2mm}%
\fcolorbox{black}{white}{\parbox{\dimexpr\linewidth-2\fboxsep-2\fboxrule\relax}{\centering\textbf{Introduction to Robotics Homework \##1 -- Draft}}}\par
\textbf{#2}\par}
\newcommand{\conventions}{Use right-handed rotations, column vectors, radians,
and meters. $\R{A}{B}$ maps coordinates from frame $B$ into frame $A$;
$\T{A}{B}$ also includes the origin of $B$ expressed in $A$.}
\newcommand{\submission}{\textbf{Submit:} labeled hand derivations and sketches,
your changed Python source files, the requested numerical/visual evidence,
and an AI-use statement (including ``none'' if applicable). Show your reasoning
before using code to check hand calculations. Public checks are intentionally
limited. The instructor supplies the deadline and submission destination.}

\begin{document}
\courseheader{2}{Rotation matrices and coordinate frames}
\textbf{Schedule:} week 2, approximately one week. \conventions

\begin{enumerate}
\item \textbf{Implement the basic rotation functions.}
In \path{src/intro_robotics/transforms.py}, complete the HW02 methods
\texttt{rot2}, \texttt{rotx}, \texttt{roty}, \texttt{rotz}, and
\texttt{rot\_inv}. The supplied input checks explain invalid arguments.
\begin{enumerate}
\item First write the $2\times2$ XY rotation and the three $3\times3$
named-axis rotations by hand. Define a positive rotation using the right-hand rule.
\item Implement those matrices directly with NumPy. Use the rotation-matrix
property for inversion; do not call \path{numpy.linalg.inv}.
The supplied \texttt{axis\_rotation} helper is for physical joints; do not
call it in these assigned functions.
\item In a course-container shell at \path{/workspace}, run:
\begin{lstlisting}
python -m homework.hw02.hw02_rotations
\end{lstlisting}
Add two checks of your own at nonzero angles. Explain what each checks.
\end{enumerate}

\item \textbf{Distinguish fixed-axis and current-axis composition.}
Start with a frame aligned with $W$.
\begin{enumerate}
\item Write the matrix product for a rotation by $a$ about X, then $b$
about Z, then $c$ about Y, with all axes belonging to the fixed world frame.
Write the product again when each rotation uses the current moving frame.
Do not expand these symbolic products.
\item Now use just two rotations: $+\pi/2$ about X followed by $+\pi/2$
about Z. Calculate both final matrices by hand: first fixed axes, then
current axes. For each, find the final local X axis expressed in $W$.
Sketch the result and check the products using your functions.
\end{enumerate}

\item \textbf{Describe a camera using its axes.}
The camera's local Z axis points along $[0,1,0]^T$ in the world. Its local Y
axis points along $[0,0,-1]^T$ in the world.
\begin{enumerate}
\item Find its local X axis using right-handed geometry, then construct
$\R{W}{C}$ by arranging the three axes as columns.
\item Map the camera-frame vector $[1,2,3]^T$ into $W$. Check that its
length is preserved. Explain what each column of $\R{W}{C}$ means.
\end{enumerate}
\end{enumerate}

\newpage
\begin{enumerate}[start=4]
\item \textbf{Find a relative rotation.}
Two frames have orientations relative to a common reference $A$:
\[
\R{A}{B}=\begin{bmatrix}1&0&0\\0&0&-1\\0&1&0\end{bmatrix},\qquad
\R{A}{C}=\begin{bmatrix}0&-1&0\\1&0&0\\0&0&1\end{bmatrix}.
\]
\begin{enumerate}
\item Derive the product for $\R{B}{C}$ and evaluate it by hand.
Use frame labels to explain why the multiplication order is correct.
\item Map $[1,0,0]^T$ from $C$ into $B$, then from $B$ into $A$.
Compare with mapping directly from $C$ into $A$.
\end{enumerate}

\item \textbf{Visualize and update an orientation.}
Run the supplied viewer inside the same ROS/Docker environment:
\begin{lstlisting}
python -m homework.hw02.hw02_rotations --show
\end{lstlisting}
The figure appears on the browser desktop from HW01; move or minimize RViz
if it hides the Matplotlib window. Red, green, and blue lines are local
X, Y, and Z. A save-only alternative is:
\begin{lstlisting}
python -m homework.hw02.hw02_rotations --save tmp/hw02.png
\end{lstlisting}
\begin{enumerate}
\item In your own short script, create a \texttt{VizScene}, add the world
frame and a second frame, and display either result from problem 2(b).
The provided \texttt{add\_frame} accepts a $4\times4$ matrix: start with
\texttt{np.eye(4)} and put your rotation into its top-left $3\times3$ block.
Keep the example's fixed drawing offset $[0.5,0,0]^T$ in the final column
so the two frames are easy to distinguish; change only the rotation block.
\item Use \texttt{scene.update(As=[T\_world, T\_moving])} to change the
second frame to the other result. Supply one transform for every frame you
added. Save both views and explain the direction of the local X axis.
\item Optional: vary one angle smoothly and animate the frame. Implementing
the drawing library is not part of the assignment.
\end{enumerate}
\end{enumerate}

\textbf{What to understand:} matrix columns are frame axes; composition
order depends on the axes used; a rotation inverse has a special structure.
No other orientation parameterizations are required.

\submission
\end{document}
